\chapter{Logiciel et reproductibilité}
\label{chap:08-logiciel}

Pineau demandait en 2018 un travail reproductible, réutilisable et robuste. Ce chapitre décrit
comment le dépôt du projet, \repo, tente de satisfaire ces trois exigences. Il présente
l'organisation du dépôt, le paquet Python qui réunit les calculs réutilisables, les formats de
fichiers, la vérification mécanique de chaque nombre publié, les tests, l'intégration continue,
les environnements figés, le matériel et le calcul consommé. Un dernier point dit ce qui manque au
dépôt et comment l'obtenir.

\section{L'organisation du dépôt}
\label{sec:depot}

Le dépôt suit un flux à sens unique, des essais vers les documents
(\cref{fig:architecture}). Un essai écrit un fichier de résultats brut : une courbe d'évaluation au
format CSV pour une simulation, un journal JSONL pour un essai GRPO, un fichier JSON de lecture pour
un article du sondage. Un script d'analyse lit ces fichiers et écrit un fichier JSON de résultats.
Deux scripts de vérification recalculent ensuite chaque nombre des textes à partir de ces fichiers
JSON. Aucun nombre d'un document ne vient d'ailleurs.

\fig{architecture}{Le dépôt, première partie : des essais aux fichiers JSON. Les entraînements écrivent les résultats bruts (CSV, JSONL, journaux). Les scripts d'analyse, appuyés sur le paquet \file{pineau\_rl}, en tirent un fichier JSON par expérience. \texttt{make analysis} refait toute cette chaîne sans aucun entraînement.}

\fig{architecture-2}{Le dépôt, seconde partie : des fichiers JSON aux documents. \file{check\_numbers.py} et \file{check\_article.py} confrontent chaque nombre des textes à son fichier JSON, \file{tools/tables.py} génère les tableaux de ce rapport. \texttt{make verify} enchaîne analyses, contrôles et tests puis vérifie que chaque JSON revient à l'identique. L'intégration continue refait ce chemin à chaque envoi sur GitHub.}

\begin{table}[tbp]
\centering
\small
\caption{Les dossiers du dépôt et leur rôle.}
\label{tab:depot}
\begin{tabularx}{\linewidth}{@{}lX@{}}
\toprule
Chemin & Contenu \\
\midrule
\file{study/} & Le sondage : protocole (\file{PROTOCOL.md}), fiche du lecteur (\file{GRILLE.md}),
  tirages, lectures (\file{codes/}, \file{codes\_v2/}, \file{codes\_x8/}), harmonisation
  (\file{overrides*.json}), provenance des lecteurs (\file{readers/}), scripts d'analyse et fichiers JSON de
  résultats, dont ceux des écarts 30 à 33 (\file{survey\_extra.py}, \file{x8\_stats.py},
  \file{q5\_versions\_stats.py}). \\
\file{sim/} & Les simulations : protocole (\file{SIM\_PROTOCOL.md}), un essai (\file{run.py}), une
  réévaluation sous physique modifiée (\file{dose.py}), listes d'essais (\file{jobs*.txt}),
  lanceurs et surveillant (\file{drive.sh}, \file{queue.sh}, \file{worker.sh}, \file{scale.sh},
  \file{watchdog.sh}), courbes (\file{results/}, avec les deux vidéos d'évaluation de E6 dans
  \file{results/e6\_video/}), journaux (\file{logs/}), politiques finales (\file{models/}), l'agent
  publié par le Zoo et ses fichiers (\file{zoo\_hub/}). \\
\file{sim/analysis/} & Un script et un fichier JSON par expérience (E1 à E7, B1, B3, dose,
  correction de Bonferroni), avec leurs variantes à vingt graines aléatoires et à budget égal. Les analyses
  des écarts 21 à 31, une par fichier : décompte strict de E2 (\file{e2\_strict.py}), renversement de
  \task{Hopper} testé (\file{e7\_test.py}), E1 sur seize couples, E3 avec Holm, E5 hors échantillon,
  calibration de la règle, facteur de B1, courbes, tailles d'effet, vitesse des agents de E6. Le relevé des
  hyperparamètres (\file{hyperparams.py}). Trois analyses écrites d'avance, en attente de leurs
  entraînements (\file{c6\_determinism.py}, \file{e6\_gamma.py}, \file{c2.py}). \\
\file{sim/llm/} & Le volet GRPO : protocole, \file{grpo.py}, lanceurs, diagnostics
  (\file{diag/}), journaux, résultats, \file{summary.py} et \file{logs\_stats.py} (écart~19). La relance
  de l'écart~20 : son lanceur \file{c7.sh} et son analyse \file{c7\_stats.py}, écrite avant tout
  essai. \\
\file{src/pineau\_rl/} & Le paquet Python des calculs réutilisables (\cref{sec:paquet}). \\
\file{tests/} & Les tests unitaires (\cref{sec:tests}). \\
\file{envs/} & Les environnements figés, la plateforme et la chaîne de documents
  (\cref{chap:D-environnements}). \\
\file{docs/} & Le document de synthèse en deux langues, sa bibliographie, ses figures et
  \file{check\_numbers.py}. \\
\file{article/} & L'article de blog en deux langues, ses figures et \file{check\_article.py}. \\
\file{report/} & Le présent rapport technique, ses tableaux générés et ses figures. \\
\file{tools/} & \file{compare\_results.py}, qui compare les fichiers JSON valeur par valeur. \\
\bottomrule
\end{tabularx}
\end{table}

\section{Le paquet \texttt{pineau\_rl}}
\label{sec:paquet}

Les calculs qui servent à plusieurs scripts forment un paquet Python installable,
\file{src/pineau\_rl/}, version 1.0.0, sous licence BSD 3-Clause. Ses seules dépendances sont
NumPy, SciPy, statsmodels et rliable. Chaque script d'analyse l'importe. Le \file{Makefile} place
\file{src/} dans le chemin de recherche de Python, si bien que les scripts tournent sans installer
le paquet. Les sous-modules se chargent à la demande : \texttt{import pineau\_rl.stats} ne charge
pas rliable. Le paquet compte quatre modules publics (\texttt{stats}, \texttt{bootstrap}, \texttt{runs}, \texttt{grid}) et un module de mise en forme (\texttt{fmt}).

\subsection{\texttt{pineau\_rl.stats}}

Ce module réunit les statistiques du sondage et des expériences. Chaque fonction est déterministe.
L'en-tête du module avertit que les calculs suivent l'ordre exact des scripts d'origine, pour que
les fichiers de résultats se régénèrent à l'octet près : simplifier une expression change le
dernier chiffre d'un résultat.
\begin{itemize}
  \item \texttt{wilson(k, n, z=1.96)} et \texttt{wilson\_ci(k, n, alpha=0.05)} : l'intervalle de
        \textcite{wilson1927probable} d'une proportion $k/n$. Avec $\hat p = k/n$, son centre vaut
        $(\hat p + z^2/2n)/(1 + z^2/n)$ et sa demi-largeur
        $z\sqrt{\hat p(1-\hat p)/n + z^2/4n^2}\,/(1 + z^2/n)$. La seconde fonction délègue à
        statsmodels, avec $z = \Phi^{-1}(1-\alpha/2)$.
  \item \texttt{newcombe\_diff(ka, na, kb, nb)} : l'intervalle à 95~\% de la différence de deux
        proportions par la méthode hybride de \textcite{newcombe1998interval}, construite à partir
        des deux intervalles de Wilson $(l_a, u_a)$ et $(l_b, u_b)$ :
        $d - \sqrt{(\hat p_a - l_a)^2 + (u_b - \hat p_b)^2}$ et
        $d + \sqrt{(u_a - \hat p_a)^2 + (\hat p_b - l_b)^2}$, avec $d = \hat p_a - \hat p_b$.
  \item \texttt{fisher(ka, na, kb, nb)} : la valeur $p$ bilatérale du test exact de Fisher sur le
        tableau $2 \times 2$.
  \item \texttt{compare(a, b)} : les deux précédentes réunies, \texttt{\{diff, lo, hi, p\}}, en
        pleine précision.
  \item \texttt{holm(pvalues)} : les valeurs $p$ ajustées par la procédure descendante de
        \textcite{holm1979simple}, $\tilde p_{(i)} = \min\bigl(1, \max_{j \le i}(m-j+1)\,p_{(j)}\bigr)$.
  \item \texttt{cochran\_armitage(k, n, x)} : le test de tendance linéaire des proportions
        \parencite{armitage1955tests}. Avec $\bar p = \sum k_i / \sum n_i$ et
        $\bar x = \sum x_i n_i / \sum n_i$, la statistique vaut
        $z = \sum x_i (k_i - n_i \bar p) \big/ \sqrt{\bar p(1-\bar p)\sum n_i (x_i - \bar x)^2}$.
  \item \texttt{cohen\_kappa(a, b)} : le kappa de \textcite{cohen1960coefficient},
        $(p_o - p_e)/(1 - p_e)$, \texttt{None} quand $p_e = 1$.
  \item \texttt{probability\_of\_improvement(x, y)} : $P(X > Y) + \tfrac12 P(X = Y)$ sur toutes les
        paires d'essais, soit la statistique $U$ de Mann-Whitney divisée par $n_x n_y$
        \parencite{agarwal2021precipice}.
  \item \texttt{welch\_curve(a, b)}, \texttt{welch\_pooled(a, b)} et
        \texttt{henderson\_script(M, S, a, b, n)} : les trois lectures de la statistique de
        \textcite{henderson2018matters} étudiées par l'expérience B1. La première applique le test
        de \textcite{welch1947generalization} à chaque point d'évaluation, puis moyenne les $t$ ou
        les $p$. La deuxième empile tous les points des courbes comme des observations
        indépendantes. La troisième refait le calcul du script \file{from\_same\_dist.py}
        \parencite{henderson2017samedist} : un test de Welch construit à partir des moyennes et des
        écarts types entre épisodes, où la dispersion entre graines aléatoires n'entre jamais.
\end{itemize}

\begin{lstlisting}[language=Python,caption={Le calcul du script de Henderson, dans \texttt{pineau\_rl.stats}.},captionpos=b]
def henderson_script(M, S, a, b, n):
    t, p = _st.ttest_ind_from_stats(M[a].mean(0), S[a].mean(0), n,
                                    M[b].mean(0), S[b].mean(0), n, equal_var=False)
    return t.mean(), p.mean()
\end{lstlisting}

\subsection{\texttt{pineau\_rl.bootstrap}}

Ce module enveloppe les intervalles bootstrap stratifiés de rliable
\parencite{agarwal2021precipice}. \texttt{iqm(x)} rend la moyenne interquartile, moyenne de la
moitié centrale des scores. \texttt{iqm\_interval(scores, reps=2000)} rend, pour chaque groupe
d'essais, la moyenne interquartile et son intervalle bootstrap par percentile à 95~\%.
\texttt{improvement\_interval(x, y, reps=2000, confidence=0.95)} rend la probabilité d'amélioration
et son intervalle, au niveau voulu : l'analyse de Bonferroni de l'écart 18 des simulations s'en sert
avec $1 - 0{,}05/24$. \texttt{probability\_of\_improvement\_bootstrap(x, y, rng)} rééchantillonne
$x$ et $y$ avec un générateur NumPy explicite, sans rliable.

Un défaut de rliable 1.2.0 commande l'usage de ce module. La bibliothèque accepte un argument
\texttt{random\_state} mais tire ses rééchantillons dans le générateur global de NumPy (écart 13 des
simulations). Un script fixe donc ce générateur une seule fois, au début
(\texttt{np.random.seed(2026)}), puis appelle les fonctions dans un ordre fixe : chaque appel
consomme le générateur global, si bien qu'un intervalle dépend des appels qui le précèdent. Le
module importe rliable dès son propre chargement, jamais dans une fonction, pour qu'aucun import ne
survienne après la graine aléatoire.

\subsection{\texttt{pineau\_rl.runs}}

Ce module définit le format d'un essai de simulation (\cref{sec:formats}) et ses lecteurs.
\texttt{run\_path(root, algo, task, seed)} donne le chemin d'un essai. \texttt{read\_run(path,
final\_step=None, upto=None)} lit un fichier et le vérifie : colonnes requises présentes, numéros de pas
strictement croissants. Elle rend \texttt{None} pour un fichier absent ou un essai inachevé. L'argument
\texttt{upto} coupe la courbe à un budget plus petit, ce qui sert au budget égal.
\texttt{write\_run(path, rows)} écrit le format depuis n'importe quelle bibliothèque.
\texttt{finals(\dots)} rend le score de chaque graine aléatoire au pas voulu, \texttt{load\_grid(\dots)} la
grille complète \texttt{\{tâche: \{algorithme: scores\}\}} avec la liste des couples manquants.

\subsection{\texttt{pineau\_rl.grid}}

Ce module lit la grille du sondage. Il définit les listes de critères (\texttt{V1} : Q1 à Q8 et X1
à X7 ; \texttt{V2} : Y1 à Y12), les critères entiers (\texttt{INT} : X1, X6, Y3), les valeurs
permises et les types d'incertitude de Y11 (\texttt{KINDS}).
\texttt{load\_reads(root, readers)} charge les fichiers de lecture,
\texttt{load\_overrides(path)} l'harmonisation. La valeur finale d'un critère suit l'une de deux
sémantiques, gardées à part parce que les scripts diffèrent : \texttt{value\_v1} (une lecture et
son harmonisation) et \texttt{value} (la lecture de la grille v2 quand elle a le critère, sinon celle
de la grille v1, avec Y3 changé en oui ou non à partir de 2). \texttt{tally(values)} rend le
numérateur et le dénominateur d'une part, les \texttt{na} exclus. \texttt{validate(rec, \dots)}
contrôle un fichier de lecture (\cref{sec:formats}).

\subsection{\texttt{pineau\_rl.fmt}}

Ce module écrit les nombres comme les textes les écrivent. \texttt{half\_up(x, nd)} arrondit la
moitié vers l'extérieur, comme on l'apprend à l'école ($1{,}15 \to 1{,}2$), alors que la fonction
\texttt{round} de Python arrondit au pair. Le flottant passe par sa représentation la plus courte,
pour que 1,15 reste 1,15. \texttt{pct}, \texttt{num}, \texttt{ci} et \texttt{pts} produisent une
part en pourcentage, un entier avec séparateur de milliers, un intervalle en points entiers et les
bornes signées d'une différence, en style français ou anglais.

\section{Les formats de fichiers}
\label{sec:formats}

\paragraph{Un essai de simulation} Un fichier CSV par programme, tâche et graine aléatoire,
\file{sim/results/<algo>/<tâche>/seed<k>.csv}, une ligne par évaluation (\cref{tab:format-essai}).
\begin{table}[htbp]
\centering
\small
\caption{Les deux premières lignes du fichier d'un essai de simulation : PPO PyTorch sur \task{HalfCheetah}, graine aléatoire 0. Sous chaque intitulé, le nom de la colonne dans le fichier CSV. La case vide est une mesure non faite.}
\label{tab:format-essai}
\begin{tabular}{@{}rrrrr@{}}
\toprule
Pas & Retour moyen & Écart type & Retour moyen, physique modifiée & Secondes écoulées \\
{\footnotesize\texttt{step}} & {\footnotesize\texttt{eval\_mean}} & {\footnotesize\texttt{eval\_std}} & {\footnotesize\texttt{perturbed\_mean}} & {\footnotesize\texttt{wall}} \\
\midrule
\num{50000} & \num{329.4} & \num{26.6} & & \num{85.5} \\
\num{100000} & \num{690.4} & \num{71.0} & & \num{188.4} \\
\bottomrule
\end{tabular}
\end{table}
\texttt{eval\_mean} et \texttt{eval\_std} donnent la moyenne et l'écart type des retours des dix
épisodes d'évaluation, \texttt{perturbed\_mean} le retour moyen sous physique modifiée (vide quand il
n'est pas mesuré), \texttt{wall} les secondes écoulées. Un essai est fini quand son dernier pas
atteint le budget du programme. Le \cref{tab:format-essai} donne les deux premières lignes de PPO
PyTorch sur \task{HalfCheetah}, graine aléatoire 0, arrondies au dixième.

\paragraph{Un essai GRPO} Un fichier JSONL par modèle, récompense et graine aléatoire
(\cref{chap:07-modeles-de-langue}). Les lignes d'évaluation contiennent \texttt{step},
\texttt{greedy}, \texttt{pass1}, \texttt{pass8} et \texttt{truncated}. Les lignes d'entraînement,
tous les dix pas, \texttt{train\_right}, \texttt{truncated}, \texttt{loss}, \texttt{grad\_norm},
\texttt{retries} et \texttt{wall}. Ces fichiers ne gardent aucun verdict par problème. La relance de
l'écart~20 ajoute, pour chaque essai, un fichier \file{seed<graine>.items.jsonl.gz} qui garde le verdict, la
longueur et la réponse de chaque problème à chaque évaluation. Son journal écrit une ligne à chaque pas.

\paragraph{Une lecture d'article} Un fichier JSON par article et par lecteur, rangé
sous \file{study/codes\_v2/} dans le dossier du lecteur. Pour chaque critère, la valeur, une citation exacte de
l'article et sa page, ou la mention \enquote{absent après lecture des sections expériences, annexes
et liens} avec une page nulle (\cref{chap:A-grille}) :
\begin{lstlisting}
{"id": "K001", "coder": "H", "title": "...",
 "criteria": {
  "Q1":  {"value": "yes", "quote": "We evaluate on six MuJoCo tasks", "page": 6},
  "X1":  {"value": 5, "quote": "averaged over 5 random seeds", "page": 7},
  "Y11": {"value": "yes", "kind": "se",
          "quote": "shaded areas denote one standard error", "page": 7}},
 "doubts": "..."}
\end{lstlisting}

\paragraph{Le validateur} \texttt{pineau\_rl.grid.validate} rend la liste des problèmes d'un fichier
de lecture, vide quand il est valide. Il vérifie le schéma : chaque critère présent, des valeurs
\texttt{yes}, \texttt{no} ou \texttt{na}, des entiers ou \texttt{null} pour X1, X6 et Y3, une famille
X5 parmi \texttt{llm}, \texttt{classic} et \texttt{theory}, un champ \texttt{kind} quand Y11 vaut oui,
une citation pour chaque critère. Il vérifie les règles de portée des écarts 3 et 14 du protocole :
\begin{itemize}
  \item Q1 non impose \texttt{na} partout ;
  \item Y2 et Y11 valent \texttt{na} si et seulement si X2 vaut non ;
  \item Y4 vaut \texttt{na} quand Y3 vaut 1 ;
  \item Y5 vaut \texttt{na} si et seulement si Q4 vaut \texttt{na}.
\end{itemize}
Avec une fonction de recherche dans le texte du PDF, le validateur vérifie enfin que chaque citation
s'y trouve. La
fonction \texttt{segments} retire d'abord les notes du lecteur entre crochets, coupe une citation
cousue par \enquote{[...]} en morceaux et ignore les morceaux de moins de douze lettres. Le script
\file{study/check\_codes.py} applique ce contrôle avec les PDF. La fiche du lecteur lui demande de
le lancer sur ses fichiers jusqu'à zéro problème.

\section{Chaque nombre vérifié contre son fichier}
\label{sec:verification}

Deux scripts vérifient les textes publiés. Ils ne lisent aucun nombre écrit à la main : ils le
recalculent.

\paragraph{\file{docs/check\_numbers.py}} Pour chaque affirmation chiffrée du document de synthèse,
le script recalcule la valeur à partir d'un fichier de résultats, la met en forme comme chaque
langue l'écrit (\texttt{pineau\_rl.fmt}), puis cherche la chaîne obtenue dans le texte, espaces et
retours à la ligne normalisés. Quelques comptes du sondage que les scripts d'analyse n'écrivent
pas sont recomptés depuis les lectures, avec les mêmes règles. Des assertions protègent les faits
qualitatifs, par exemple le nom de la tâche où le programme en tête change sous physique modifiée. Le script
vérifie 188 affirmations dans les deux langues et s'arrête au premier nombre manquant :
\file{docs/build.sh} ne compile donc pas un document dont un chiffre a dérivé.

\begin{lstlisting}[language=Python,caption={Une affirmation de \file{docs/check\_numbers.py} : la chaîne attendue est construite depuis les données, jamais tapée.},captionpos=b]
add("agreement",
    f"{same} décisions sur {total} ({pct(same / total)})",
    f"{same} decisions out of {total} ({pct(same / total, lang='en')})")
\end{lstlisting}

\paragraph{\file{article/check\_article.py}} Le script de l'article de blog ajoute un contrôle de
couverture. Chaque affirmation construit, depuis un fichier de résultats, la phrase que les deux
langues écrivent, puis la cherche dans le texte. Les quelques nombres sans fichier (une
diapositive, un article cité, une taille d'échantillon) sont cherchés dans le document de
synthèse lui-même. Le script retire ensuite du texte toutes les phrases vérifiées et les noms
propres qui contiennent des chiffres : tout nombre qui reste est un échec. Les lignes de la
bibliographie sont confrontées au fichier \file{docs/refs.bib}. L'article compte 161 affirmations
vérifiées dans chaque langue.

\section{Les cibles du \file{Makefile}}
\label{sec:make}

Aucune cible n'entraîne quoi que ce soit. L'entraînement passe par \file{sim/run.py},
\file{sim/drive.sh} et \file{sim/llm/grpo.py}.

\begin{description}[style=nextline,leftmargin=1.5em]
  \item[\texttt{make install}] installe le paquet en mode éditable dans l'environnement d'analyse,
        sans dépendance téléchargée : elles viennent de \file{envs/analysis.lock}.
  \item[\texttt{make analysis}] relance tous les scripts d'analyse, dans l'ordre des dépendances, par trois
        cibles qui se lancent aussi seules.
        \texttt{analysis-study} lance les douze scripts du sondage, dont ceux des écarts 30 à 33.
        \texttt{analysis-sim} lance les dix-sept scripts des simulations à dix graines aléatoires, dont le
        décompte strict de E2 et le test du renversement de \task{Hopper}, puis l'analyse statistique du
        pipeline officiel du Zoo, les variantes des écarts 12 et 20 (vingt graines aléatoires, réplication de
        E4 sur les graines aléatoires 10 à 19, budget égal) et les huit analyses des écarts 22 à 31, qui lisent
        les fichiers des précédentes. \texttt{analysis-llm} lance \file{summary.py} à trois et à cinq graines
        aléatoires, puis \file{logs\_stats.py} (écart~19). Les variantes passent par des variables
        d'environnement (\texttt{SEEDS}, \texttt{SUFFIX}, \texttt{BUDGET}). Deux analyses de l'écart~33,
        \file{e6\_enjoy.py} et \file{e6\_hub.py}, demandent l'environnement du Zoo et les politiques :
        aucune cible ne les relance, comme \file{e6\_official\_eval.json}.
  \item[\texttt{make hyperparams}] relit les réglages des cinq programmes dans les bibliothèques de
        l'environnement d'entraînement et écrit \file{sim/analysis/hyperparams.json} (écart~28). Cette cible
        reste hors de \texttt{verify}, qui tourne dans l'environnement d'analyse.
  \item[\texttt{make analysis-suite}] lance les trois analyses écrites d'avance pour les expériences pas
        encore finies (écarts 32 à 34). Leurs fichiers JSON n'entrent dans la vérification qu'une fois
        leurs entraînements finis.
  \item[\texttt{make check}] lance les deux scripts de vérification de la \cref{sec:verification}.
  \item[\texttt{make test}] lance les tests unitaires avec pytest.
  \item[\texttt{make lint}] lance ruff et mypy sur le paquet et les tests.
  \item[\texttt{make verify}] enchaîne \texttt{analysis}, \texttt{check} et \texttt{test}, puis
        demande à git si chaque fichier JSON de résultats est revenu à l'octet près.
  \item[\texttt{make verify-head}] fait la même chose dans un export neuf du dernier commit, sans
        les politiques, les essais du Zoo ni les journaux d'entraînement, qu'aucune analyse ne lit.
        Cette cible laisse intacte une copie de travail où des essais écrivent encore.
  \item[\texttt{make figures}] régénère les figures du document de synthèse
        (\file{docs/figures/build.py}).
  \item[\texttt{make docs}] lance \texttt{check}, puis construit le document de synthèse
        (\file{docs/build.sh}).
  \item[\texttt{make report}] construit le présent rapport (\file{report/build.sh}) : tableaux
        générés depuis les fichiers JSON, figures, puis LuaLaTeX et biber, en français et en
        anglais.
\end{description}

\begin{lstlisting}[language=bash,caption={Recalculer tous les nombres depuis les résultats commités, sans entraînement.},captionpos=b]
# Créer l'environnement d'analyse
python3.11 -m venv sim/env-an
# Installer les paquets de l'analyse et des tests
sim/env-an/bin/pip install -r envs/analysis.txt -r envs/test.txt
# Analyses, nombres des textes, tests, diff octet par octet
make verify
# La même chose dans un export neuf du dernier commit
make verify-head
\end{lstlisting}

\section{L'intégration continue}
\label{sec:ci}

Le fichier \file{.github/workflows/verify.yml} rejoue la vérification à chaque envoi et à chaque
demande de fusion, sur deux plateformes. Les actions GitHub employées sont fixées par empreinte de
commit, pas par étiquette. Chaque tâche installe l'environnement d'analyse depuis
\file{envs/analysis.lock}, chaque paquet contrôlé par son empreinte SHA-256
(\texttt{pip install -{}-require-hashes}), puis installe le paquet, lance les tests, relance toutes
les analyses, vérifie les nombres des textes et compare chaque fichier de résultats au commit avec
\file{tools/compare\_results.py}.
\begin{itemize}
  \item Sur Ubuntu 24.04 (x86-64), avec Python 3.11.18, la version que nous employons, la
        comparaison accepte un écart relatif de $10^{-12}$ sur les flottants. Elle exige l'égalité
        des comptes, des verdicts, des classements et de toutes les chaînes. Les valeurs $p$ bougent à leur
        quinzième ou seizième chiffre avec la bibliothèque mathématique.
  \item Sur macOS 15 (arm64), système de la famille de machines dont viennent les résultats (un Mac
        sous macOS 27.2, \cref{sec:log-calcul}), la comparaison exige l'égalité exacte.
\end{itemize}
Une troisième tâche construit l'image du \file{Dockerfile} et y lance les tests et la vérification
des nombres. L'image part de \texttt{python:3.11-slim}, fixée par empreinte, installe le même
fichier verrouillé et monte le dépôt au lancement au lieu de le copier : la même image sert à
n'importe quel commit.

\section{Les tests}
\label{sec:tests}

Le dossier \file{tests/} contient huit fichiers de tests et un fichier d'outils communs
(\file{conftest.py}). Les tests importent le paquet et la couche mince
\file{sim/analysis/runs.py}, jamais un script d'analyse, qui calcule dès son import. Trois sortes de
contrôles s'y côtoient : des valeurs publiées dans la littérature, un calcul indépendant
(statsmodels, SciPy, NumPy) et des valeurs de référence sur des données fabriquées, figées à
$10^{-12}$ près en relatif. Une valeur de référence qui change annonce un fichier de résultats qui
change.

\begin{description}[style=nextline,leftmargin=1.5em]
  \item[\file{test\_stats.py}] Wilson contre une valeur publiée, contre statsmodels et sur ses
        propriétés (l'intervalle contient la part et reste dans $[0,1]$, il est symétrique, il s'élargit à 99~\%) ; Newcombe contre une valeur publiée et
        reconstruit depuis les deux intervalles de Wilson ; Fisher contre SciPy ; Holm calculé à la
        main et ses propriétés ; Cochran-Armitage égal à $\sqrt N$ fois la corrélation de Pearson
        entre score et issue ; le kappa de Cohen et sa liste de catégories ; la probabilité
        d'amélioration.
  \item[\file{test\_bootstrap.py}] la moyenne interquartile ; les intervalles de rliable, identiques
        sous une graine aléatoire globale fixée et dépendants de cette seule graine aléatoire ; la probabilité
        d'amélioration bootstrap, qui n'emploie que son propre générateur.
  \item[\file{test\_henderson.py}] les trois lectures de la statistique de Henderson contre SciPy ;
        la valeur $p$ du $t$ moyen avec les degrés de liberté de Welch moyens ; le calcul du script,
        qui ignore la dispersion entre graines aléatoires : sur six graines aléatoires fabriquées, il atteint un $t$ de 9
        là où un test sur les graines aléatoires ne voit rien.
  \item[\file{test\_runs.py}] le format d'un essai : écriture puis lecture, essais absents ou
        inachevés rendus \texttt{None}, fichiers mal formés qui lèvent une erreur, scores finaux et
        grille de référence.
  \item[\file{test\_grid.py}] les deux sémantiques de valeur, la conversion de Y3, l'harmonisation,
        le validateur (fichier valide, chaque problème signalé, règles de portée, citations) et le
        découpage des citations.
  \item[\file{test\_formatting.py}] l'arrondi de la moitié vers l'extérieur, les pourcentages, les
        milliers, les intervalles et le signe moins des deux langues.
  \item[\file{test\_pipelines.py}] la lecture des essais par \file{sim/analysis/runs.py} et le
        résumé GRPO de bout en bout, sur trois graines aléatoires fabriquées : un essai inachevé est écarté, le
        gain moyen et l'intervalle bootstrap sont ceux attendus.
  \item[\file{test\_draw.py}] le tirage de 2026 revient de sa règle (\file{study/draw\_2026.py}) et
        dépend de sa graine aléatoire.
\end{description}

\section{Les environnements figés}
\label{sec:envs}

Quatre environnements Python ont fait tourner le projet, tous sous Python 3.11.18. Chaque fichier de
\file{envs/} est le \texttt{pip freeze} de l'un d'eux, avec la version de
Python et la plateforme en tête. Les fichiers \file{*.lock} ajoutent l'empreinte SHA-256 de chaque
paquet. Le \cref{chap:D-environnements} donne les versions. L'environnement d'entraînement
(\file{envs/train.txt}) sert à \file{run.py} et \file{dose.py}, l'environnement d'analyse
(\file{analysis.txt}) aux analyses, à la vérification et à l'intégration continue, l'environnement
du Zoo (\file{zoo.txt}) au pipeline officiel de l'expérience E6, l'environnement des modèles de
langue (\file{llm.txt}) à GRPO et aux figures du document de synthèse. Un cinquième fichier,
\file{test.txt}, ajoute pytest à l'environnement d'analyse.

\section{Les politiques entraînées}
\label{sec:release}

Les politiques finales permettent de réévaluer un agent sans le réentraîner, ce qu'a fait la
réévaluation par degrés dans des physiques modifiées (\file{sim/dose.py}). Celles de PPO, de TRPO et de
l'expérience E6 sont dans le dépôt. Les poids de SAC et de CrossQ, trop lourds pour git, sont
attachés à la publication GitHub \texttt{policies-v1}, sous forme d'archives tar listées avec leur
empreinte SHA-256 dans un fichier \file{SHA256SUMS}. Le script \file{sim/fetch\_models.py}
télécharge chaque archive, vérifie son empreinte, la décompresse sous \file{models/}, puis vérifie
chaque fichier de politique contre \file{sim/models/MANIFEST.sha256}, les empreintes prises
sur la machine qui les a entraînées. Les téléchargements se suivent un par un,
séparés par une pause aléatoire de 1,5 à 4 secondes. Une requête refusée est reprise après une
attente tirée au hasard entre zéro et un plafond qui double à chaque échec (recul exponentiel à gigue
complète), cinq tentatives au plus.

\begin{lstlisting}[language=bash,caption={Récupérer et vérifier les politiques de SAC et de CrossQ.},captionpos=b]
# Télécharger, décompresser, tout vérifier
cd sim && python3 fetch_models.py
# Vérifier seulement les politiques présentes
cd sim && python3 fetch_models.py --check
\end{lstlisting}

\section{Le matériel et le calcul}
\label{sec:log-calcul}

Tous les essais ont tourné sur un seul ordinateur portable, un Mac à puce Apple M2 Max sous macOS
27.2 (arm64), partagé avec d'autres projets. Chaque entraînement de simulation tourne sur un seul
fil d'exécution. Le nombre de tâches en parallèle a changé au fil du projet : quatre cœurs au départ,
jusqu'à huit accordés chaque heure selon la charge de la machine (écart 15 des simulations), puis
deux cœurs pour tout le projet (écart 19). Ce partage ne touche pas les
résultats tant qu'un essai ne dépend que de sa graine aléatoire : les relances le vérifient pour PPO,
l'écart 32 le teste encore pour SAC et CrossQ (\cref{sec:registres}). La \cref{sec:calcul} décrit les lanceurs,
le surveillant et la file d'attente. Le \cref{tab:calcul} donne, par programme et par tâche, le
nombre d'essais, le nombre de pas, la durée médiane d'un essai et les heures de cœur consommées.

Les essais GRPO ont des durées très différentes selon l'appareil. D'après le champ \texttt{wall} de
leurs journaux au pas 100, un essai Qwen a pris de 42 à 52 minutes sur GPU et de 4,3 à 11,4 heures
sur CPU, un essai Llama de 69 à 148 minutes sur GPU. Ces écarts reflètent aussi la charge variable
de la machine. La relance de l'écart~20 prévoit, d'après ces journaux, environ 58 heures de GPU pour ses
quarante essais, à lancer machine au repos.

\section{Ce qui n'est pas dans le dépôt}
\label{sec:absent}

Les PDF des articles lus sont protégés par le droit d'auteur. Le script \file{study/fetch.py} les
télécharge de nouveau depuis les sources publiques (\texttt{cd study \&\& python3 fetch.py
final\_2026.json pdf}). Les diapositives de Pineau ne sont pas dans le dépôt pour la même raison :
elles se trouvent sur les sites d'ICLR 2018 et de NeurIPS 2018. Le fichier
\file{slides-en.md} les décrit une à une. Les journaux des sessions de travail avec les agents qui
ont lu les articles ne sont pas publics. La consigne exacte, le modèle et les dates de chaque
lecteur, reconstruits à partir de ces journaux, le sont (\file{study/readers/}).

\section{Licence et citation}
\label{sec:licence}

Le code, les protocoles, les données dérivées et les documents sont publiés sous la licence
BSD 3-Clause, celle de scikit-learn (fichier \file{LICENSE}, Copyright (c) 2026, Warith Harchaoui).
Le fichier \file{CITATION.cff} donne la référence, que GitHub propose en APA et en BibTeX.
